\documentclass{article}
\usepackage[utf8]{inputenc}
\usepackage[margin=2cm,includefoot]{geometry}
\usepackage{booktabs}
\usepackage{graphicx}
\usepackage{amsmath, amssymb, amsthm}
\usepackage{physics}
\usepackage[shortlabels]{enumitem}
\usepackage{cancel}
\usepackage{lastpage}
\usepackage{braket}
\usepackage{mathrsfs}
\usepackage{bm}

\usepackage{fancyhdr}
\pagestyle{fancy}
\newcommand{\AssignmentNum}{1}
\lhead{Stromingsleer en transportverschijnselen (NS-265B)\\
2025/2026}
\rhead{Tereza Bílková (7137257), Bela F. Brunner (8002193),\\ Niels Epema (2927578), Adam Zich (0074187)}
\fancyfoot{}
\fancyfoot[R]{Page \thepage /\pageref{LastPage}}

\newcommand{\vba}{\mathbf{A}}
\newcommand{\vbb}{\mathbf{B}}
\newcommand{\C}{\mathbf{C}}
\newcommand{\x}{\mathbf{\hat{x}}}
\newcommand{\y}{\mathbf{\hat{y}}}
\newcommand{\z}{\mathbf{\hat{z}}}
\newcommand{\vr}{\mathbf{\hat r}}
\newcommand{\thet}{\boldsymbol{{\hat \theta}}}
\newcommand{\ph}{\boldsymbol{{\hat \phi}}}
\newcommand{\ba}{\mathbf{a}}
\newcommand{\bb}{\mathbf{b}}
\newcommand{\vl}{\mathbf{l}}
\newcommand{\vv}{\mathbf{v}}

\newcommand{\R}{\mathbb{R}}
\newcommand{\p}{\mathbb{P}}
\newcommand{\N}{\mathbb{N}}
\newcommand{\Z}{\mathbb{Z}}
\newcommand{\Q}{\mathbb{Q}}
\newcommand{\I}{\mathbb{I}}
\newcommand{\E}{\mathbb{E}}

\newcommand{\mbeq}{\overset{!}{=}}

\begin{document}
    \begin{center}
    \LARGE{\bfseries{ Project \AssignmentNum}}
    \line(1,0){430}
    \end{center}

\subsection*{c)}
First, we show that $df/dz=\frac{df/dZ}{dZ/dz}$ by invoking the chain rule: If $y=f(u)$ and $u=g(x)$, then $\frac{dy}{dx}=\frac{dy}{du}\cdot \frac{du}{dx}$. Since $f$ is a function of $Z$ and $Z$ can be expressed as a function of $z$. Accordingly:
\begin{align*}
    \frac{df}{dz}=\frac{df}{dZ}\cdot \frac{dZ}{dz}=\frac{df/dZ}{dz/dZ}\quad \square 
\end{align*}
We can now evaluate both of the term separately. We first look at the bottom term. For this, we can just plug in the formula for $z(Z)=Z+\dfrac{1}{Z}$, where we used that $b=1$. Plugging this in gives: 
\begin{align*}
    \frac{dz}{dZ}=(Z+\frac{1}{Z})'=1-\frac{1}{Z^2}
\end{align*}
Secondly, we are looking at the top term: 
\begin{align*}
    \frac{df}{dZ}&=\left( U e^{-i\alpha} Z + \frac{U e^{i\alpha}}{Z}- \frac{i\Gamma}{2\pi} \ln(Z)\right)'\\
&=U e^{-i\alpha} - \frac{U e^{i\alpha}}{Z^2}- \frac{i\Gamma}{2\pi Z} 
\end{align*}
Now putting both of them together: 
\begin{align*}
    \frac{df}{dz}&=\frac{U e^{-i\alpha} - \frac{U e^{i\alpha}}{Z^2}- \frac{i\Gamma}{2\pi Z} }{1-\frac{1}{Z^2}}\\
    &=\frac{ZU e^{-i\alpha} - \frac{U e^{i\alpha}}{Z}- \frac{i\Gamma}{2\pi} }{Z-\frac{1}{Z}}
\end{align*}
Now, we replace $Z=e^{i\gamma}$, where we set the radius to be 1: 
\begin{align*}
    \frac{df}{dz}&=\frac{Ue^{i\gamma} e^{-i\alpha} - {U e^{i\alpha}}e^{-i\gamma}- \frac{i\Gamma}{2\pi} }{e^{i\gamma}-e^{-i\gamma}}\\
    &=\frac{U(e^{i(\gamma-\alpha)}-e^{-i(\gamma-\alpha)})-\frac{i\Gamma}{2\pi} }{{e^{i\gamma}-e^{-i\gamma}}}\\
    &=\frac{U2i\sin(\gamma-\alpha )-\frac{i\Gamma}{2\pi}}{2i\sin (\gamma)}\\
    &=\frac{U\sin(\gamma -\alpha)-\Gamma/4\pi}{\sin(\gamma )} \quad \square 
\end{align*}

\subsection*{d)}
We first want to find a relationship to relate $z$ to gamma. For this we can use the definition, $b=1$ and $R=1$: 
\begin{align*}
    z=Z+\frac{1}{Z}=e^{i\gamma }+e^{-i\gamma }=2\cos(\gamma)
\end{align*}
With this formula, we can find for $z=\pm 2$:
\begin{align*}
    2=2\cos(\gamma )\Leftrightarrow \gamma_{2z} =\arccos (1)=0 \\
    -2=2\cos(\gamma )\Leftrightarrow \gamma_{-2z} =\arccos (-1)=\pi 
\end{align*}
If we know plug this into the formula obtained in c), we get: 

\begin{align*}
    \frac{df(2)}{dz}
=
\frac{U \sin(0 - \alpha) - \Gamma/(4\pi)}
{\sin(0)}=\frac{U \sin(-\alpha) - \Gamma/(4\pi)}
{0}\\
\frac{df(-2)}{dz}
=
\frac{U \sin(\pi - \alpha) - \Gamma/(4\pi)}
{\sin(\pi)}=\frac{U \sin(\pi-\alpha) - \Gamma/(4\pi)}
{0}
\end{align*}
Therefore, in both cases the velocity actually approaches infinity when it gets closer to the edges, since we are dividing by 0. 

\subsection*{e)}
The correct choice for $\Gamma$ would be such that the upper part of the fraction also becomes zero, since then we could apply L'Hospital and find a real velocity. Since, as given by the problem, we are only concerned with the trailing edge, we only have to look at $z=2$, as $0\leq\alpha \leq \frac{\pi}{2}$ and the wind is standardly coming from the left, therefore the leading edge will always be at $z=-2$. Finding $\Gamma$:
\begin{align*}
    U\sin(-\alpha)-\Gamma /(4\pi)=0\\
    \Leftrightarrow \Gamma =-{U\sin(\alpha)}{4\pi}
\end{align*}
We can now calculate the velocity at the edge with L'Hospital: 
\begin{align*}
    \frac{df(2)}{dz}=\lim _{\gamma \to 0 }U\frac{\sin(\gamma -\alpha)-\sin(-\alpha)}{\sin(\gamma)}=\lim _{\gamma \to 0 }U\frac{\cos(\gamma -\alpha)}{\cos(\gamma)}=U\frac{\cos(-\alpha)}{1}=U\cos (-\alpha)
\end{align*}
Since the Blasius theorem is just the foundation for the theorem of Kutta-Joukowski, we can just use this one. With this, we get: 
\begin{align*}
    \mathcal F_x-i\mathcal F_y&=i\rho U\Gamma \\
    &=-i\rho U^2{\sin(\alpha)}{4\pi}
\end{align*}
Since we defined $U$ to be real, this gives: 
\begin{align*}
        \mathcal F_x-i\mathcal F_y&=-i\rho U^2{\sin(\alpha)}{4\pi} \\
    \Leftrightarrow \mathcal F_y&=4\pi \rho U^2\sin (\alpha )\\
    \Rightarrow \mathcal F_y&\propto \sin(\alpha)
\end{align*}
Since $\sin (0)=0$ and $\sin(\pi/2)=1$, we can see that the lift force is 0 if the wind comes directly from the side and becomes maximum the wind comes directly from the bottom, which makes sense intuitively. 

We can now calculate the units: 
\begin{align*}
    [\mathcal F_y]=[ \rho U^2b]=\frac{\text{kg}\cdot \text{m}^2}{\text{m}^3\cdot \text{s}^2}\cdot\text{m}=\frac{\text{kg}}{ \text{s}^2}=\frac{\text {N}}{\text{m}}
\end{align*}
The formula was derived with the non-dimensionalized length of $b$. We have to restore it at this point, so that we arrive at Newton per meter. This is also what is to be expected, as we are looking at a two-dimensional problem and the wind can only attack a one dimensional line instead of an area, which would be the real-world equivalent.
\end{document}
